% Extra coding exercise: bias-variance trade-off.
% Paste into the report after \section*{Chapter 3}. Requires booktabs, graphicx, amsmath, float
% and \graphicspath{{figures/}} in the preamble.
% Figures: bv_mse.png, bv_knn_decomp.png, bv_lm.png, generated by code/export_bv_figures.py
% (seed 0). Copy code/figures/ to the Overleaf project.

\section*{Extra exercise: bias--variance trade-off}

\subsection*{Setup}

The ground truth is $y = f(x) = \sin(2x_1)$. A training set of $n_{\text{train}} = 50$ points and a test set of $n_{\text{test}} = 500$ points are drawn once, with all $p = 20$ predictors independent and uniform on $[-1, 1]$. Only $x_1$ enters $f$; the other 19 predictors are pure noise. Training responses are $y_i = f(x_i) + \varepsilon_i$ with $\varepsilon_i \sim \mathcal N(0, 0.5^2)$, while the test responses are kept noise-free so that the test error decomposes exactly into squared bias and variance.

For each number of predictors $p \in \{1, 2, 3, 4, 10, 20\}$ the models are fitted on the first $p$ columns of the design matrix. The training noise is redrawn $n_{\text{rep}} = 100$ times; each repetition yields a linear regression fit and nine KNN fits with $K = 1, \dots, 9$, and the predictions on the fixed test set are stored as $\hat f^{(l)}(x_r)$, $r = 1, \dots, 500$, $l = 1, \dots, 100$.

\subsection*{Estimating bias and variance}

For every test point the average prediction over the repetitions, $\bar f(x_r) = \frac{1}{100}\sum_l \hat f^{(l)}(x_r)$, estimates $\mathbb E[\hat f(x_r)]$. The three quantities are then averaged over the test points:
\begin{align*}
  \text{Bias}^2 &= \frac{1}{n_{\text{test}}} \sum_r \bigl(f(x_r) - \bar f(x_r)\bigr)^2, \\
  \text{Var}    &= \frac{1}{n_{\text{test}}\, n_{\text{rep}}} \sum_r \sum_l \bigl(\hat f^{(l)}(x_r) - \bar f(x_r)\bigr)^2, \\
  \text{MSE}    &= \frac{1}{n_{\text{test}}\, n_{\text{rep}}} \sum_r \sum_l \bigl(\hat f^{(l)}(x_r) - f(x_r)\bigr)^2 .
\end{align*}
Because the test responses contain no noise, $\text{MSE} = \text{Bias}^2 + \text{Var}$ holds exactly; in the simulation the largest discrepancy is $10^{-16}$. In the code, with \texttt{yhat} of shape $(500, 100)$ and \texttt{ytest} of length 500, the three lines are
\begin{verbatim}
ybar     = yhat.mean(axis=1)
biasSQ   = np.mean((ytest - ybar)**2)
variance = np.mean((yhat - ybar[:, None])**2)
err      = np.mean((yhat - ytest[:, None])**2)
\end{verbatim}
where \texttt{[:, None]} turns the length-500 vector into a column so that it is broadcast across the 100 repetitions.

\subsection*{Results}

\begin{figure}[H]
  \centering
  \includegraphics[width=\linewidth]{bv_mse.png}
  \caption{Test MSE of KNN regression as a function of $1/K$ (log scale; $1/K = 1$ is $K = 1$), one panel per number of predictors $p$. The dashed line is the test MSE of linear regression with the same $p$. Compare ISLR Figure~3.20.}
  \label{fig:bv-mse}
\end{figure}

\begin{figure}[H]
  \centering
  \includegraphics[width=\linewidth]{bv_knn_decomp.png}
  \caption{Decomposition of the KNN test MSE into squared bias and variance as a function of $1/K$, one panel per $p$.}
  \label{fig:bv-knn}
\end{figure}

\begin{figure}[H]
  \centering
  \includegraphics[width=0.6\linewidth]{bv_lm.png}
  \caption{Decomposition of the linear regression test MSE as a function of the number of predictors $p$.}
  \label{fig:bv-lm}
\end{figure}

\begin{table}[H]
  \centering
  \begin{tabular}{rrrrrrr}
    \toprule
    & \multicolumn{3}{c}{Linear regression} & \multicolumn{3}{c}{KNN at its best $K$} \\
    \cmidrule(lr){2-4} \cmidrule(lr){5-7}
    $p$ & Bias$^2$ & Var & MSE & $K^*$ & MSE & Bias$^2$ / Var \\
    \midrule
    1  & 0.028 & 0.010 & 0.038 & 9 & 0.032 & 0.004 / 0.028 \\
    2  & 0.027 & 0.015 & 0.042 & 9 & 0.045 & 0.020 / 0.025 \\
    3  & 0.028 & 0.018 & 0.046 & 8 & 0.080 & 0.050 / 0.031 \\
    4  & 0.032 & 0.026 & 0.058 & 6 & 0.141 & 0.098 / 0.043 \\
    10 & 0.036 & 0.070 & 0.107 & 8 & 0.339 & 0.306 / 0.032 \\
    20 & 0.054 & 0.178 & 0.233 & 9 & 0.384 & 0.357 / 0.028 \\
    \bottomrule
  \end{tabular}
  \caption{Squared bias, variance and test MSE for linear regression, and for KNN at the value of $K$ that minimises its MSE, for each $p$ (seed 0, 100 repetitions).}
  \label{tab:bv}
\end{table}

\subsection*{Discussion}

\paragraph{Linear regression.}
Its squared bias is almost constant, between $0.027$ and $0.054$ (Figure~\ref{fig:bv-lm}). A straight line cannot represent $\sin(2x_1)$ on $[-1, 1]$, and adding noise predictors does not change that. Its variance, however, grows from $0.010$ at $p = 1$ to $0.178$ at $p = 20$, because every additional predictor is one more coefficient to estimate from the same 50 noisy observations, and the noise coefficients wander around zero from sample to sample. The MSE of the linear model is therefore driven by variance as $p$ grows.

\paragraph{KNN and the choice of $K$.}
Within each panel of Figure~\ref{fig:bv-knn} the variance falls monotonically with $K$, from about $0.25$ at $K = 1$ to $0.03$ at $K = 9$, in close agreement with $\sigma^2/K = 0.25/K$: the prediction is the average of $K$ noisy responses. This variance curve is essentially the same for every $p$. The squared bias moves the other way, rising with $K$ once the neighbourhood becomes wide enough to average over the curvature of $\sin(2x_1)$, and it rises much more steeply when $p$ is large. The optimal $K$ is where the two curves cross: for $p = 1$ and $2$ the bias is negligible and the largest $K$ tried is best, whereas for $p = 4$ the minimum lies at $K = 6$ with a visible U-shape in the MSE.

\paragraph{The curse of dimensionality.}
The dominant effect of $p$ on KNN is through the bias. With $p = 1$ the nearest neighbours of a test point are genuinely close in $x_1$, so their average is close to $f$; the bias is $0.004$. With $p = 20$ the Euclidean distance is dominated by the 19 noise coordinates, so the ``nearest'' neighbours are spread over the whole $x_1$ range and their average is far from $f(x_1)$; the bias rises to $0.36$ and no choice of $K$ repairs it. This is why the KNN curves in Figure~\ref{fig:bv-mse} shift upward with $p$ while their shape stays the same.

\paragraph{KNN versus linear regression.}
For $p = 1$ and $2$, KNN with a large $K$ beats linear regression (MSE $0.032$ vs.\ $0.038$ at $p = 1$), because the true relationship is non-linear and the flexible method can follow it at low cost in variance. From $p = 3$ onwards the linear model wins for every $K$, and at $p = 20$ its MSE is $0.23$ against $0.38$ for the best KNN. The linear model pays for each noise predictor with a modest increase in variance; KNN pays with a large increase in bias. This reproduces the pattern of ISLR Figure~3.20: a parametric method with the wrong functional form can still outperform a non-parametric one when the number of predictors is moderate and the sample is small, because the non-parametric method's ability to adapt is exactly what makes it vulnerable to irrelevant dimensions.
